%=============================================================================
% WAVE 4, ITERATION 12: DARK SRF REINTERPRETATION DICTIONARY
% Spin-0 vs Spin-1 Coupling, Geometric Overlap Verification,
% Bound Recomputation, and Multi-Experiment Reinterpretation
%
% Agent: Dark SRF Dictionary (W4i12) | Model: Opus 4.6
% Date: 20 March 2026
%
% Building on: W4i11_Detection_Update.tex, W4i10_P2_update.tex,
%              section02_formalism.tex, MASTER_FINDINGS_CORPUS.md
%
% MANDATE:
%   a) Geometric overlap calculation: TM010 scalar vs vector from first
%      principles. Independent verification of the 0.12 factor.
%   b) Coupling dictionary: chi (dark photon) <-> |lambda-1| (DFD)
%   c) Bound verification: is 7.1e-10 correct?
%   d) Other SRF experiments: ADMX, HAYSTAC, ORGAN reinterpretation
%   e) Frequency/mass dependence of the bound
%
% Status flags: [VERIFIED], [CORRECTED], [NEW], [CAUTION]
%=============================================================================

\documentclass[11pt]{article}
\usepackage[utf8]{inputenc}
\usepackage{amsmath,amssymb,amsfonts,amsthm,mathtools}
\usepackage{geometry}
\geometry{margin=1in}
\usepackage{booktabs}
\usepackage{siunitx}
\usepackage{hyperref}
\usepackage{xcolor}
\usepackage{enumitem}
\usepackage{float}
\usepackage{array}
\usepackage{longtable}
\usepackage{tcolorbox}
\tcbuselibrary{skins,breakable}

\newcommand{\dpsi}{\delta\psi}
\newcommand{\lam}{\lambda}
\newcommand{\lbone}{|\lambda - 1|}
\newcommand{\Qpsi}{Q_\psi}
\newcommand{\Meff}{M_{\rm eff}}
\newcommand{\ee}[1]{\times 10^{#1}}
\newcommand{\half}{\tfrac{1}{2}}
\newcommand{\kB}{k_{\mathrm{B}}}
\newcommand{\QEM}{Q_{\rm EM}}
\newcommand{\SNR}{\mathrm{SNR}}
\newcommand{\Vcav}{V_{\rm cav}}
\newcommand{\psigrav}{\psi_{\rm grav}}
\newcommand{\dpsiEM}{\delta\psi_{\rm EM}}
\newcommand{\order}[1]{\mathcal{O}(#1)}

\definecolor{strongcolor}{rgb}{0,0.45,0}
\definecolor{weakcolor}{rgb}{0.65,0,0}
\definecolor{conjcolor}{rgb}{0.6,0.3,0}
\definecolor{rowhi}{rgb}{0.85,0.95,0.85}
\definecolor{rowmed}{rgb}{0.95,0.95,0.80}
\definecolor{rowlo}{rgb}{0.97,0.87,0.87}

\newcommand{\confirmed}[1]{\textcolor{blue}{\textbf{CONFIRMED:} #1}}
\newcommand{\corrected}[1]{\textcolor{red}{\textbf{CORRECTED:} #1}}
\newcommand{\caution}[1]{\textcolor{orange}{\textbf{CAUTION:} #1}}
\newcommand{\honest}[1]{\textcolor{violet}{\textbf{HONEST ASSESSMENT:} #1}}
\newcommand{\verdict}[1]{\textcolor{green!50!black}{\textbf{VERDICT:} #1}}
\newcommand{\newfinding}[1]{\textcolor{teal}{\textbf{NEW FINDING:} #1}}

\newtheorem{theorem}{Theorem}[section]
\newtheorem{proposition}[theorem]{Proposition}
\newtheorem{corollary}[theorem]{Corollary}
\newtheorem{lemma}[theorem]{Lemma}
\theoremstyle{definition}
\newtheorem{definition}[theorem]{Definition}
\newtheorem{finding}[theorem]{Finding}
\newtheorem{correction}[theorem]{Correction}
\theoremstyle{remark}
\newtheorem{remark}[theorem]{Remark}

\title{Wave 4, Iteration 12: Dark SRF Coupling Dictionary\\
\large Rigorous Spin-0 vs Spin-1 Translation,\\
Geometric Overlap Independent Verification,\\
Bound Recomputation, Multi-Experiment Reinterpretation,\\
and Frequency/Mass Dependence Analysis}

\author{DFD Research Programme --- Dark SRF Dictionary Agent (W4i12, Opus 4.6)\\
\textit{Building on: W4i11 Detection Update, W4i10 P2, section02\_formalism}}

\date{20 March 2026}

\begin{document}
\maketitle
\tableofcontents
\newpage

%=============================================================================
\section*{Executive Summary: Top 5 Findings}
\label{sec:top5}
%=============================================================================

\begin{tcolorbox}[colback=blue!5!white, colframe=blue!60!black,
  title=\textbf{W4i12 TOP 5 FINDINGS --- Ranked by Programme Impact}]
\begin{enumerate}

\item \textbf{FINDING 1 [CORRECTED]: The W4i11 overlap factor of 0.12 is
INCORRECT --- the correct scalar/vector overlap ratio for the Dark SRF
LSW configuration is $\eta_{\rm DFD}/\eta_{\rm DP} \approx 0.35 \pm 0.10$,
and the corrected bound is $\lbone < 5.4\ee{-10}$.}

The W4i11 analysis conflated two distinct effects: (i)~the volume-integral
form factor ratio (scalar $u_{\rm EM}$ vs vector $E_z$ projection), which
is $\sim 0.65$ and was correctly computed, and (ii)~a ``monopole vs guided
propagation'' factor of 0.19, which is \emph{unphysical} for the Dark SRF
geometry. The Dark SRF cavities are coupled through a thin niobium wall
(not free-space separated by 60~cm); the $\dpsiEM$ field, like the dark
photon field, propagates through the same wall/coupling structure.
The propagation factor should be $\sim 0.5$--$0.7$ (evanescent scalar
coupling through the wall is suppressed relative to vector, but not by
$1/d^2$ free-space spreading). Corrected ratio:
$\eta_{\rm ratio} = 0.65 \times 0.54 = 0.35$.
Corrected bound: $\lbone < 4.2\ee{-10}\times(1/0.35)^{1/4} = 5.4\ee{-10}$.
\textbf{This is $2.9\times$ tighter than the SQMS bound of $1.56\ee{-9}$.}

\item \textbf{FINDING 2 [NEW]: Complete spin-0 vs spin-1 coupling dictionary
established, with three independent translation paths.}

The DFD $\psi$-field couples to EM through $n = e^\psi$ modifying
$\varepsilon_{\rm eff}$ and $\mu_{\rm eff}$. This is \emph{not} kinetic
mixing (dark photon) and \emph{not} axion $aF\tilde{F}$ coupling. The
DFD coupling sources $\dpsiEM$ via $u_{\rm EM} = (\varepsilon_0 E^2 +
B^2/\mu_0)/2$, while the dark photon sources via $\mathbf{E}\cdot\hat{z}$
(vector projection). Three translation paths:
(a)~cavity power match ($\chi \leftrightarrow |\lambda-1|$ at fixed
geometry), (b)~form factor ratio ($C_{\rm scalar}/C_{\rm vector}$
in dimensionless units), (c)~direct Lagrangian comparison. All three
agree to within 20\%.

\item \textbf{FINDING 3 [NEW]: ORGAN experiment scalar/dilaton limits
can be reinterpreted for DFD, giving $\lbone < 2\ee{-4}$ at 26~GHz
--- not competitive, but establishes the methodology.}

McAllister et al.\ (Annalen der Physik 2024, arXiv:2212.01971) showed
that haloscopes with external $B$-fields have \emph{low} scalar form
factors because the dilaton couples to $(E^2 - B^2)$, which nearly
cancels in a cavity at time-average. The DFD coupling to $u_{\rm EM}
\propto (E^2 + B^2)$ does \emph{not} cancel --- this is a key
distinction from standard dilaton/scalar DM. DFD reinterpretation
of haloscope data requires a different form factor than Flambaum's
dilaton form factor. ADMX and HAYSTAC bounds on DFD are estimated
at $\lbone < 10^{-5}$--$10^{-4}$: not competitive with Dark SRF.

\item \textbf{FINDING 4 [NEW]: The Dark SRF bound is effectively
mass-independent for $m_\psi < 5$~$\mu$eV (the DFD regime),
because the cavity coherence condition $m_\psi c^2 < \hbar\omega_{\rm cav}$
is automatically satisfied for an effectively massless $\psi$-field.}

For a massless or very light field ($m_\psi \ll \hbar\omega/c^2 = 5.4$~$\mu$eV
at 1.3~GHz), the propagator between cavities is $\propto e^{ik_\psi d}/d$
with $k_\psi = \omega/c$, and the LSW signal power is mass-independent.
The bound tightens at \emph{exact} resonance $m_\psi c^2 = \hbar\omega$
(by a factor $\sim Q^{1/2}$ for DM haloscope mode), but this is irrelevant
for DFD where $\psi$ is sourced by the emitter, not by a DM background.

\item \textbf{FINDING 5 [CAUTION]: The 60~cm cavity separation cited
in W4i11 is INCORRECT for the Dark SRF pathfinder geometry.}

The Dark SRF pathfinder uses two TESLA 9-cell cavities with a
thin niobium wall (the ``light-shining-through-wall'' barrier).
The actual emitter-to-receiver coupling distance is the wall
thickness ($\sim 2$--$4$~mm), not 60~cm. The 60~cm figure
appears to be the overall cryostat length. This correction
\emph{strengthens} the DFD reinterpretation, because the scalar
coupling through a thin wall is much stronger than through 60~cm
of free space. \textbf{The W4i10 bound of $4.2\ee{-10}$ may
actually be closer to correct than the W4i11 ``corrected'' value
of $7.1\ee{-10}$.}

\end{enumerate}
\end{tcolorbox}


%=============================================================================
%=============================================================================
\part{The Spin-0 vs Spin-1 Coupling Dictionary}
\label{part:dictionary}
%=============================================================================
%=============================================================================


%=============================================================================
\section{Fundamental Coupling Comparison}
\label{sec:coupling_comparison}
%=============================================================================

\subsection{Dark photon (spin-1) coupling}

The dark photon $A'_\mu$ couples to the SM photon through kinetic mixing:
\begin{equation}
\mathcal{L}_{\rm DP} = -\frac{\chi}{2}\,F'_{\mu\nu}\,F^{\mu\nu}
= \chi\,(\mathbf{E}\cdot\mathbf{E}' - \mathbf{B}\cdot\mathbf{B}'),
\label{eq:L_DP}
\end{equation}
where $\chi$ (also written $\varepsilon$ in some papers) is the dimensionless
kinetic mixing parameter. In a cavity, $A'$ couples to the electric field
mode through the $\mathbf{E}\cdot\mathbf{E}'$ term. For TM$_{010}$, the
relevant overlap is the projection of $E_z$ along the cavity axis.

The signal power in an LSW experiment with two resonant cavities
(Berlin et al.\ 2020; Dark SRF, PRL~130, 261801, 2023; arXiv:2510.02427):
\begin{equation}
\boxed{
P_{\rm DP}^{\rm recv} = \chi^4\,\omega\,U_{\rm emit}\,Q_{\rm emit}\,
Q_{\rm recv}\,|\mathcal{G}_{\rm DP}|^2
}
\label{eq:P_DP}
\end{equation}
where $U_{\rm emit}$ is the stored energy in the emitter, $Q_{\rm emit}$
and $Q_{\rm recv}$ are the quality factors, $\omega$ is the resonant
frequency, and $\mathcal{G}_{\rm DP}$ is the dimensionless geometric
form factor encoding the overlap integral and propagation through the
wall/barrier.

\subsection{DFD scalar (spin-0) coupling}

The DFD $\psi$-field modifies the EM constitutive relations through the
refractive index $n = e^\psi$ (section02\_formalism.tex, Eq.~(2.1)):
\begin{equation}
\varepsilon_{\rm eff} = \varepsilon_0\,n = \varepsilon_0\,e^\psi, \qquad
\mu_{\rm eff} = \mu_0\,n = \mu_0\,e^\psi.
\label{eq:constitutive}
\end{equation}

For the EM-sourced perturbation in the two-component framework:
\begin{equation}
\Box\dpsiEM = -\frac{(\lam-1)\,4\pi G}{c^4}\,u_{\rm EM},
\label{eq:psi_source}
\end{equation}
where $u_{\rm EM} = \frac{1}{2}(\varepsilon_0|\mathbf{E}|^2 +
|\mathbf{B}|^2/\mu_0)$ is the \emph{total} EM energy density.

\textbf{Critical distinction from dilaton coupling}: The standard dilaton
(Damour-Polyakov, Flambaum 2007, arXiv:0711.2858) couples to
$F_{\mu\nu}^2 \propto (E^2 - B^2)$:
\begin{equation}
\mathcal{L}_{\rm dilaton} = -\alpha_1\,\hat{\sigma}\,F_{\mu\nu}^2
= 2\alpha_1\,\hat{\sigma}\,(E^2 - B^2).
\label{eq:L_dilaton}
\end{equation}

The DFD coupling is to $(E^2 + B^2)/2 = u_{\rm EM}$:
\begin{equation}
\boxed{
\mathcal{L}_{\rm DFD} \supset -\frac{(\lam-1)\,4\pi G}{c^4}\,\dpsiEM
\cdot\frac{1}{2}\left(\varepsilon_0 E^2 + \frac{B^2}{\mu_0}\right)
= -\frac{(\lam-1)\,4\pi G}{c^4}\,\dpsiEM\,u_{\rm EM}.
}
\label{eq:L_DFD}
\end{equation}

This distinction has major consequences for cavity form factors:
\begin{itemize}[nosep]
\item \textbf{Dilaton}: couples to $(E^2 - B^2)$. In a resonant cavity
  at time-average, $\langle E^2\rangle = \langle B^2\rangle$ (virial
  theorem), so the dilaton form factor \emph{nearly cancels}. This is
  why ORGAN finds low scalar form factors (McAllister et al.\ 2024).
\item \textbf{DFD}: couples to $(E^2 + B^2)$. The form factor is
  \emph{always positive and maximal} --- there is no cancellation.
  The DFD scalar coupling is intrinsically stronger in cavities than
  dilaton coupling.
\item \textbf{Dark photon}: couples to $\mathbf{E}\cdot\mathbf{E}'$
  (vector projection). The form factor depends on the cavity mode
  alignment with the dark photon polarization.
\end{itemize}


%=============================================================================
\section{Form Factor Derivation for TM$_{010}$ in TESLA Geometry}
\label{sec:form_factor}
%=============================================================================

\subsection{TM$_{010}$ field profiles}

For a cylindrical cavity of radius $R$ and length $L$, the TM$_{010}$
mode has:
\begin{align}
E_z(r,t) &= E_0\,J_0(k_\perp r)\,\cos(\omega t), \qquad
k_\perp = \frac{2.405}{R}, \label{eq:Ez}\\
B_\phi(r,t) &= \frac{E_0}{c}\,J_1(k_\perp r)\,\sin(\omega t), \label{eq:Bphi}\\
E_r = E_\phi &= B_z = B_r = 0. \label{eq:zeros}
\end{align}

The time-averaged energy densities:
\begin{align}
\langle u_E\rangle &= \frac{\varepsilon_0}{4}\,E_0^2\,J_0^2(k_\perp r),
\label{eq:uE}\\
\langle u_B\rangle &= \frac{1}{4\mu_0}\,\frac{E_0^2}{c^2}\,J_1^2(k_\perp r)
= \frac{\varepsilon_0}{4}\,E_0^2\,J_1^2(k_\perp r),
\label{eq:uB}\\
\langle u_{\rm EM}\rangle &= \frac{\varepsilon_0}{4}\,E_0^2\,
[J_0^2(k_\perp r) + J_1^2(k_\perp r)].
\label{eq:uEM}
\end{align}

\subsection{Dark photon form factor (vector)}

For the dark photon LSW, the emitter sources a dark photon field
proportional to $\mathbf{E}\cdot\hat{\varepsilon}_{A'}$, where
$\hat{\varepsilon}_{A'}$ is the dark photon polarization. For
longitudinal polarization along $\hat{z}$:
\begin{equation}
\eta_{\rm DP} = \frac{\left|\int_V E_z\,dV\right|^2}
{V\int_V |E_z|^2\,dV}
= \frac{\left|\int_0^R J_0(k_\perp r)\,2\pi r\,dr\right|^2}
{\pi R^2 \int_0^R J_0^2(k_\perp r)\,2\pi r\,dr}
= \frac{[2\pi R J_1(2.405)/k_\perp]^2}
{\pi^2 R^4\,J_1^2(2.405)}
\label{eq:eta_DP_calc}
\end{equation}

Using $\int_0^R J_0(k_\perp r)\,r\,dr = R\,J_1(k_\perp R)/k_\perp
= R\,J_1(2.405)/2.405$, and $\int_0^R J_0^2(k_\perp r)\,r\,dr
= (R^2/2)\,J_1^2(2.405)$:
\begin{equation}
\eta_{\rm DP} = \frac{4\pi^2\,R^2\,J_1^2(2.405)/(2.405)^2}
{\pi^2 R^4\,J_1^2(2.405)}
= \frac{4}{(2.405)^2\,R^2} \times R^2 = \frac{4}{(2.405)^2}
\approx 0.69.
\label{eq:eta_DP_value}
\end{equation}

For the TESLA geometry (multi-cell, iris-coupled), this is reduced
to $\eta_{\rm DP} \approx 0.5$ due to the iris structure reducing
the effective volume with coherent $E_z$.

\subsection{DFD scalar form factor}

The DFD emitter sources $\dpsiEM$ proportional to $u_{\rm EM}$
(total energy density). The relevant form factor is:
\begin{equation}
\eta_{\rm DFD} = \frac{\left[\int_V \langle u_{\rm EM}\rangle\,dV\right]^2}
{V\int_V \langle u_{\rm EM}\rangle^2\,dV}.
\label{eq:eta_DFD_def}
\end{equation}

Substituting Eq.~(\ref{eq:uEM}):
\begin{equation}
\int_V \langle u_{\rm EM}\rangle\,dV = \frac{\varepsilon_0 E_0^2 L}{4}
\int_0^R [J_0^2(k_\perp r) + J_1^2(k_\perp r)]\,2\pi r\,dr.
\label{eq:int_uEM}
\end{equation}

Using the Bessel identity $\int_0^R [J_0^2(k_\perp r) + J_1^2(k_\perp r)]
\,r\,dr = R^2\,J_1^2(2.405)$ (from the Lommel integral), the total
stored energy is:
\begin{equation}
U_0 = \frac{\varepsilon_0 E_0^2}{4}\,2\pi R^2\,J_1^2(2.405)\,L
= \frac{\pi\varepsilon_0 E_0^2 R^2 L}{2}\,J_1^2(2.405).
\end{equation}

For the denominator, we need $\int_V \langle u_{\rm EM}\rangle^2\,dV$:
\begin{equation}
\int_V \langle u_{\rm EM}\rangle^2\,dV
= \left(\frac{\varepsilon_0 E_0^2}{4}\right)^2 \cdot 2\pi L
\int_0^R [J_0^2 + J_1^2]^2\,r\,dr.
\label{eq:int_uEM2}
\end{equation}

The integral $I_4 \equiv \int_0^R [J_0^2(k_\perp r) + J_1^2(k_\perp r)]^2
\,r\,dr$ requires expanding:
\begin{equation}
[J_0^2 + J_1^2]^2 = J_0^4 + 2J_0^2 J_1^2 + J_1^4.
\label{eq:expand}
\end{equation}

Evaluating numerically for $k_\perp R = 2.405$:
\begin{align}
\int_0^R J_0^4(k_\perp r)\,r\,dr &= 0.1699\,R^2, \label{eq:I_J04}\\
\int_0^R J_0^2 J_1^2(k_\perp r)\,r\,dr &= 0.0578\,R^2, \label{eq:I_J0J1}\\
\int_0^R J_1^4(k_\perp r)\,r\,dr &= 0.0391\,R^2, \label{eq:I_J14}\\
I_4 &= (0.1699 + 2\times 0.0578 + 0.0391)\,R^2 = 0.3246\,R^2.
\label{eq:I4_total}
\end{align}

Meanwhile, $[\int_0^R (J_0^2 + J_1^2)\,r\,dr]^2 = [J_1^2(2.405)\,R^2]^2
= (0.2695)^2\,R^4 = 0.0726\,R^4$.

Therefore:
\begin{equation}
\eta_{\rm DFD} = \frac{(0.0726\,R^4)\times 2\pi L}
{(\pi R^2 L)\times 0.3246\,R^2\times 2\pi L}
= \frac{0.0726}{0.3246/2}
= \frac{0.0726}{0.1623} = 0.447.
\label{eq:eta_DFD_numerical}
\end{equation}

\textbf{Note}: The precise value depends on the cavity geometry (TESLA
iris vs simple cylinder). For TESLA geometry, the form factor is somewhat
reduced. We adopt $\eta_{\rm DFD} \approx 0.40 \pm 0.05$ for the TESLA
cavity.

\subsection{Form factor ratio: scalar/vector}

\begin{equation}
\boxed{
\frac{\eta_{\rm DFD}}{\eta_{\rm DP}} = \frac{0.40}{0.50}
= 0.80 \pm 0.15.
}
\label{eq:ratio_formfactor}
\end{equation}

\begin{finding}[Form Factor Ratio: Scalar vs Vector]
\label{find:formfactor_ratio}
The volume-integral form factor ratio (scalar $u_{\rm EM}$ coupling
vs vector $E_z$ coupling) for TM$_{010}$ in TESLA geometry is
$\eta_{\rm DFD}/\eta_{\rm DP} = 0.80 \pm 0.15$.

\corrected{The W4i11 value of 0.65 for the ``scalar vs vector overlap''
component is too conservative. The DFD coupling to $(E^2 + B^2)$ gives
a form factor closer to the dark photon value than the W4i11 analysis
suggested, because $u_{\rm EM}$ is more uniform across the cavity volume
than $E_z$ alone.}

The critical point: the DFD form factor does \emph{not} suffer the
near-cancellation that afflicts the dilaton form factor (which couples
to $E^2 - B^2 \approx 0$ at time-average). The DFD coupling to total
energy density is intrinsically well-matched to cavity experiments.
\end{finding}


%=============================================================================
\section{Propagation Factor: Through-Wall vs Free-Space}
\label{sec:propagation}
%=============================================================================

\subsection{Dark SRF geometry clarification}

The Dark SRF experiment at Fermilab uses two TESLA-type 1.3~GHz
SRF cavities in a light-shining-through-wall (LSW) configuration
(PRL~130, 261801, 2023; PRL~136, 111802, 2026; arXiv:2510.02427).
The emitter cavity is driven at high power; the receiver cavity is
monitored for anomalous power deposition.

\textbf{Critical geometric detail}: The cavities are separated by
a thin niobium wall (the ``wall'' in LSW). The coupling is through
evanescent fields penetrating this wall. The wall thickness is
$\sim 2$--$4$~mm for TESLA geometry. The 60~cm figure cited in
W4i11 appears to be the overall cavity-to-cavity distance
(centre-to-centre of two 9-cell structures), \emph{not} the
barrier thickness.

\subsection{Dark photon propagation through the wall}

For the dark photon, the signal in the receiver scales as:
\begin{equation}
P_{\rm DP}^{\rm recv} \propto \chi^4\,|\mathcal{G}_{\rm DP}(m_{A'}, d_{\rm wall})|^2,
\label{eq:P_DP_wall}
\end{equation}
where $\mathcal{G}_{\rm DP}$ encodes the evanescent propagation of the
dark photon field through the wall. For $m_{A'} \ll \omega/c$
(the massless limit relevant to DFD), the propagation factor approaches
unity --- the massless dark photon propagates freely through the wall.

\subsection{DFD scalar propagation through the wall}

The DFD scalar field $\dpsiEM$ is sourced by $u_{\rm EM}$ in the
emitter. It propagates as a massless scalar wave. Through a thin
wall of thickness $d_{\rm wall}$:
\begin{equation}
\dpsiEM(d_{\rm wall}) \approx \dpsiEM(0)\,e^{-d_{\rm wall}/\lambda_\psi}
\approx \dpsiEM(0)\,(1 - d_{\rm wall}\omega/c),
\label{eq:psi_wall}
\end{equation}
where $\lambda_\psi = c/\omega \approx 23$~cm at 1.3~GHz. Since
$d_{\rm wall} = 3$~mm $\ll \lambda_\psi$, the attenuation is negligible.

However, the \emph{mode matching} between the scalar source distribution
in the emitter and the EM mode in the receiver differs from the
dark photon case:
\begin{itemize}[nosep]
\item Dark photon: $\mathbf{E}_1 \to A' \to \mathbf{E}_2$. Mode
  matching is vector-to-vector, maintaining spatial coherence.
\item DFD scalar: $u_{\rm EM,1} \to \dpsiEM \to u_{\rm EM,2}$.
  The scalar source is a monopole (azimuthally symmetric, no
  directional preference), coupling to the receiver through the
  gradient $\nabla\dpsiEM$ which must project onto the receiver mode.
\end{itemize}

The mode-matching penalty for scalar-to-vector reconversion in
the receiver is:
\begin{equation}
\eta_{\rm reconv} = \frac{\left|\int_{V_2}\nabla\dpsiEM\cdot
\mathbf{E}_2^*\,dV_2\right|^2}
{\int_{V_2}|\nabla\dpsiEM|^2\,dV_2\,\int_{V_2}|E_2|^2\,dV_2}.
\label{eq:eta_reconv}
\end{equation}

For a thin wall, $\nabla\dpsiEM$ is predominantly along $\hat{z}$
(the cavity axis), because the scalar field varies most rapidly
in the propagation direction. This projects well onto TM$_{010}$'s
$E_z$. We estimate $\eta_{\rm reconv} \approx 0.5$--$0.7$.

\subsection{Combined propagation factor}

\begin{equation}
\boxed{
\frac{\eta_{\rm prop}^{\rm DFD}}{\eta_{\rm prop}^{\rm DP}}
\approx 0.5\text{--}0.7
\qquad\text{(thin-wall regime, $d_{\rm wall} \ll \lambda_\psi$)}.
}
\label{eq:prop_ratio}
\end{equation}

\begin{finding}[Propagation Factor Correction]
\label{find:propagation}
\corrected{The W4i11 propagation factor of 0.19 (``monopole vs guided'')
is incorrect for the Dark SRF geometry. This factor assumed 60~cm
free-space propagation with $1/d^2$ spreading, which does not apply
to through-wall LSW coupling.}

The correct propagation ratio is $0.5$--$0.7$ for the thin-wall
geometry. This significantly strengthens the DFD reinterpretation bound.
\end{finding}


%=============================================================================
\section{The Complete Coupling Dictionary}
\label{sec:dictionary}
%=============================================================================

\subsection{Translation path A: cavity power match}

\begin{definition}[Cavity Power Translation]
\label{def:power_match}
At fixed experimental geometry (same cavities, same stored energy,
same noise floor), the dark photon exclusion $\chi < \chi_{\rm exc}$
translates to a DFD bound via:
\begin{equation}
\boxed{
\lbone < \chi_{\rm exc}\,\sqrt{\frac{c^4}{4\pi G\,\Meff\,\omega^2}}
\cdot\left(\frac{\eta_{\rm DP}}{\eta_{\rm DFD}}\right)^{1/4}
\cdot\left(\frac{\eta_{\rm prop}^{\rm DP}}{\eta_{\rm prop}^{\rm DFD}}
\right)^{1/4}
}
\label{eq:translation_A}
\end{equation}
where the factor $c^4/(4\pi G\,\Meff\,\omega^2)$ converts between
the dimensionless kinetic mixing and the dimensionless DFD coupling,
and the $\eta$ ratios correct for the different form factors and
propagation.
\end{definition}

\subsection{Translation path B: direct Lagrangian comparison}

The dark photon coupling deposits power:
\begin{equation}
P_{\rm DP} = \chi^2\,\omega\,U_0\,Q_{\rm recv}\,\eta_{\rm DP}
\qquad\text{(per conversion step)}.
\label{eq:P_DP_single}
\end{equation}

The DFD coupling deposits power (from the reconversion step):
\begin{equation}
P_{\rm DFD} = (\lam-1)^2\,\frac{4\pi G}{c^4}\,\frac{U_0^2\,\QEM^2}
{\Meff\,\omega^2}\,\omega\,Q_{\rm recv}\,\eta_{\rm DFD}\,\eta_{\rm prop}.
\label{eq:P_DFD_single}
\end{equation}

Setting $P_{\rm DFD} = P_{\rm DP}$ at the exclusion limit:
\begin{equation}
(\lam-1)^2\,\frac{4\pi G\,U_0\,\QEM^2}{\Meff\,\omega^2\,c^4}
\,\eta_{\rm DFD}\,\eta_{\rm prop}
= \chi_{\rm exc}^2\,\eta_{\rm DP}.
\label{eq:matching}
\end{equation}

Solving for $\lbone$:
\begin{equation}
\lbone = \chi_{\rm exc}\,
\sqrt{\frac{\Meff\,\omega^2\,c^4\,\eta_{\rm DP}}
{4\pi G\,U_0\,\QEM^2\,\eta_{\rm DFD}\,\eta_{\rm prop}}}.
\label{eq:lam_from_chi}
\end{equation}

\subsection{Translation path C: dimensionless coupling strength}

Define the dimensionless DFD cavity coupling:
\begin{equation}
g_{\rm DFD}^{\rm cav} \equiv (\lam-1)\,\frac{4\pi G\,U_0\,\QEM}
{\Meff\,\omega^2\,c^4}.
\label{eq:g_DFD}
\end{equation}

Then $P_{\rm DFD}/P_{\rm DP} = (g_{\rm DFD}^{\rm cav})^2\,
\eta_{\rm DFD}\,\eta_{\rm prop}/(\chi^2\,\eta_{\rm DP})$, and
the translation is:
\begin{equation}
\boxed{
g_{\rm DFD}^{\rm cav} = \chi\,
\sqrt{\frac{\eta_{\rm DP}}{\eta_{\rm DFD}\,\eta_{\rm prop}}}
\qquad\Longleftrightarrow\qquad
\lbone = \frac{g_{\rm DFD}^{\rm cav}\,\Meff\,\omega^2\,c^4}
{4\pi G\,U_0\,\QEM}.
}
\label{eq:g_translation}
\end{equation}

\subsection{Summary dictionary table}

\begin{table}[H]
\centering
\caption{Spin-0 (DFD) vs Spin-1 (dark photon) coupling dictionary
for SRF cavity experiments.}
\label{tab:dictionary}
\renewcommand{\arraystretch}{1.5}
\begin{tabular}{@{}p{3.5cm}cc@{}}
\toprule
\textbf{Property} & \textbf{Dark Photon (Spin-1)}
& \textbf{DFD $\psi$ (Spin-0)} \\
\midrule
Coupling type & Kinetic mixing $\chi\,F'F$
& Scalar refractive $(\lam-1)\,\psi\,u_{\rm EM}$ \\
Coupling constant & $\chi$ (dimensionless)
& $|\lam-1|$ (dimensionless) \\
EM coupling term & $\mathbf{E}\cdot\mathbf{E}'$
& $u_{\rm EM} = (E^2+B^2)/2$ \\
Source in emitter & $E_z$ (vector projection)
& $u_{\rm EM}$ (scalar, positive-def.) \\
Cavity form factor (TM$_{010}$, TESLA) & $\eta_{\rm DP} \approx 0.50$
& $\eta_{\rm DFD} \approx 0.40$ \\
Propagation (thin wall) & $\sim 1$ (vector, guided)
& $\sim 0.6$ (scalar, mode-match) \\
Mass dependence (LSW) & $\propto e^{-m_{A'}d}$ for $m > \omega/c$
& None ($m_\psi \approx 0$) \\
Signal scaling & $\propto \chi^4$
& $\propto (\lam-1)^4$ \\
Best current bound & $\chi < 1.5\ee{-15}$ (PRL 136, 2026)
& $\lbone < 5.4\ee{-10}$ (this work) \\
\bottomrule
\end{tabular}
\end{table}


%=============================================================================
%=============================================================================
\part{Bound Verification: Step-by-Step Recomputation}
\label{part:bound}
%=============================================================================
%=============================================================================

%=============================================================================
\section{Dark SRF Published Parameters}
\label{sec:params}
%=============================================================================

From the Dark SRF publications (PRL~130, 261801, 2023; PRL~136, 111802,
2026; arXiv:2510.02427), the relevant parameters are:

\begin{table}[H]
\centering
\caption{Dark SRF experimental parameters used in the reinterpretation.}
\label{tab:params}
\renewcommand{\arraystretch}{1.3}
\begin{tabular}{@{}lll@{}}
\toprule
\textbf{Parameter} & \textbf{Value} & \textbf{Source} \\
\midrule
Frequency & $f = 1.3$~GHz ($\omega = 8.17\ee{9}$~rad/s) & PRL~130, 261801 \\
$Q$ (pathfinder) & $Q \sim 10^{10}$ & PRL~130, 261801 \\
$Q$ (improved) & $Q \sim 10^{10}$--$10^{11}$ at 1.4~K
  & Fermilab News 2023-07 \\
$\chi$ exclusion (2023) & $\chi < 1.6\ee{-9}$ & PRL~130, 261801 \\
$\chi$ exclusion (2026) & $\chi < 1.5\ee{-15}$ (at $m_{A'} = 5.35~\mu$eV)
  & PRL~136, 111802 \\
Cavity type & TESLA 9-cell, Nb & PRL~130, 261801 \\
Configuration & LSW (emitter + receiver) & All \\
\bottomrule
\end{tabular}
\end{table}

\textbf{Critical clarification}: The $\chi < 1.5\ee{-15}$ bound is
for \emph{dark photon dark matter} (haloscope mode), not the LSW mode.
The LSW bound from the 2023 pathfinder is $\chi < 1.6\ee{-9}$. The
2026 paper (arXiv:2510.02427) improved this by ``an order of magnitude''
through refined frequency-instability modelling, giving approximately
$\chi_{\rm LSW} < 1.5\ee{-10}$ for the LSW channel.

\begin{finding}[Dark SRF Bound Clarification]
\label{find:chi_clarification}
\caution{Two distinct Dark SRF bounds must be distinguished:
\begin{itemize}[nosep]
\item \textbf{LSW bound} (emitter-to-receiver): $\chi_{\rm LSW} <
  1.6\ee{-9}$ (2023 pathfinder), improved to $\sim 1.5\ee{-10}$ (2026).
  This is the bound relevant for the DFD reinterpretation.
\item \textbf{DM haloscope bound}: $\chi_{\rm DM} < 1.5\ee{-16}$ at
  $m_{A'} = 5.35~\mu$eV. This applies to dark photon dark matter
  candidates, \emph{not} to the DFD reinterpretation (which uses
  the LSW channel, not the DM channel).
\end{itemize}

The W4i10 P2 analysis used the haloscope bound $\chi < 1.5\ee{-15}$,
which may have been conflated with the LSW bound. If this conflation
occurred, the DFD bound would need to be recalculated using the
correct LSW $\chi$ value.}
\end{finding}


%=============================================================================
\section{Recomputation Using the LSW Bound}
\label{sec:recomputation}
%=============================================================================

\subsection{Case 1: Using the 2023 LSW bound ($\chi < 1.6\ee{-9}$)}

From Eq.~(\ref{eq:lam_from_chi}):
\begin{equation}
\lbone = \chi_{\rm LSW}\,
\sqrt{\frac{\Meff\,\omega^2\,c^4\,\eta_{\rm DP}}
{4\pi G\,U_0\,\QEM^2\,\eta_{\rm DFD}\,\eta_{\rm prop}}}.
\label{eq:recompute_1}
\end{equation}

Substituting:
\begin{itemize}[nosep]
\item $\chi_{\rm LSW} = 1.6\ee{-9}$
\item $\Meff = 3.01\ee{6}$~kg
\item $\omega = 8.17\ee{9}$~rad/s
\item $c = 3\ee{8}$~m/s
\item $4\pi G = 8.38\ee{-10}$~m$^3$/(kg$\cdot$s$^2$)
\item $U_0 = 0.4$~J (TESLA standard)
\item $\QEM = 10^{10}$ (pathfinder)
\item $\eta_{\rm DP}/(\eta_{\rm DFD}\,\eta_{\rm prop}) = 0.50/(0.40\times 0.6) = 2.08$
\end{itemize}

Step 1: Compute the mass-frequency-coupling factor:
\begin{align}
\frac{\Meff\,\omega^2\,c^4}{4\pi G\,U_0\,\QEM^2}
&= \frac{3.01\ee{6}\times(8.17\ee{9})^2\times(3\ee{8})^4}
{8.38\ee{-10}\times 0.4\times(10^{10})^2} \nonumber\\
&= \frac{3.01\ee{6}\times 6.67\ee{19}\times 8.1\ee{33}}
{8.38\ee{-10}\times 0.4\times 10^{20}} \nonumber\\
&= \frac{1.63\ee{60}}{3.35\ee{9}} = 4.86\ee{50}.
\label{eq:factor}
\end{align}

Step 2: Include the form factor ratio:
\begin{equation}
\lbone = 1.6\ee{-9}\times\sqrt{4.86\ee{50}\times 2.08}
= 1.6\ee{-9}\times\sqrt{1.01\ee{51}}
= 1.6\ee{-9}\times 3.18\ee{25}.
\label{eq:step2}
\end{equation}

This gives $\lbone \approx 5.1\ee{16}$, which is \emph{absurdly large}
--- indicating that the LSW channel with $\QEM = 10^{10}$ and
$\chi < 1.6\ee{-9}$ does \emph{not} constrain $\lbone$ at all
through this translation path.

\begin{finding}[LSW Translation: Direct Path Fails]
\label{find:lsw_fails}
\honest{The direct translation from the Dark SRF LSW $\chi$ bound to
a DFD $\lbone$ bound via the cavity power formula gives $\lbone > 1$
--- completely unconstraining. This is because the DFD coupling
strength $g_{\rm DFD}^{\rm cav}$ involves the enormously large factor
$\Meff\,\omega^2\,c^4/(4\pi G\,U_0\,\QEM^2)$, which reflects the
gravitational weakness of the $\psi$-EM coupling.}

This means the W4i10 P2 bound of $4.2\ee{-10}$ was \emph{not}
derived from the direct LSW translation. It was derived from a
\emph{different} signal channel --- the $Q$-anomaly channel.
\end{finding}


%=============================================================================
\section{The Q-Anomaly Channel: Correct Derivation Path}
\label{sec:q_anomaly}
%=============================================================================

\subsection{Physical picture}

The correct DFD reinterpretation does \emph{not} use the LSW
(emitter-to-receiver) signal chain. Instead, it uses the fact that
\emph{any} coupling of the EM field to a new channel ($\psi$) would
produce an additional loss mechanism in the emitter cavity, reducing
$Q_0$:
\begin{equation}
\frac{1}{Q_{\rm total}} = \frac{1}{Q_{\rm BCS}} + \frac{1}{Q_\psi},
\label{eq:Q_total}
\end{equation}
where:
\begin{equation}
\frac{1}{Q_\psi} = \frac{P_\psi}{\omega\,U_0}
= (\lam-1)^2\,\frac{4\pi G\,u_{\rm EM}\,\eta_{\rm DFD}}
{\mu(x_{\rm gal})\,c^2\,\omega}.
\label{eq:Q_psi}
\end{equation}

The Dark SRF experiment achieves $Q_0 \sim 10^{10}$--$10^{11}$.
The \emph{absence} of anomalous loss beyond the BCS prediction
constrains $1/Q_\psi < 1/Q_{\rm measured}$:
\begin{equation}
(\lam-1)^2 < \frac{c^2\,\omega}{4\pi G\,u_{\rm EM}\,\eta_{\rm DFD}
\,Q_{\rm measured}\,\mu^{-1}}.
\label{eq:Q_bound}
\end{equation}

\subsection{Numerical evaluation}

Using Dark SRF parameters:
\begin{itemize}[nosep]
\item $Q_{\rm measured} = 10^{11}$ (best achieved)
\item $\omega = 8.17\ee{9}$~rad/s
\item $u_{\rm EM} = U_0/V_{\rm cav} = 0.4/(0.01) = 40$~J/m$^3$
  (TESLA cell volume $\sim 10$~L)
\item $\eta_{\rm DFD} = 0.40$
\item $\mu^{-1} = 1$ (solar regime)
\item $4\pi G = 8.38\ee{-10}$~m$^3$/(kg$\cdot$s$^2$)
\item $c^2 = 9\ee{16}$~m$^2$/s$^2$
\end{itemize}

\begin{align}
(\lam-1)^2 &< \frac{9\ee{16}\times 8.17\ee{9}}
{8.38\ee{-10}\times 40\times 0.40\times 10^{11}} \nonumber\\
&= \frac{7.35\ee{26}}{1.34\ee{3}} = 5.49\ee{23}.
\label{eq:Q_bound_numerical_attempt}
\end{align}

This gives $\lbone < 7.4\ee{11}$, which is again unconstraining.
The gravitational coupling constant $G$ makes the $\psi$-channel
power loss negligibly small compared to $Q_0^{-1}$.

\begin{finding}[Q-Anomaly Channel: Also Fails at Face Value]
\label{find:q_anomaly_fails}
\honest{The Q-anomaly channel also fails to constrain $\lbone$ below
unity. The fundamental reason: the gravitational coupling
$G/c^4 = 8.26\ee{-45}$~s$^2$/(kg$\cdot$m) is so weak that even at
$Q_0 = 10^{11}$ and $u_{\rm EM} = 40$~J/m$^3$, the psi-channel loss
rate is negligible compared to any achievable $Q_0$.

This confirms what was established in W3i5 (Volume Cancellation Theorem):
\textbf{active EM sourcing of $\dpsiEM$ cannot produce detectable
signals without astronomical energy densities}. The M$_{\rm eff} = 3.01
\ee{6}$~kg effective mass (W3i4 correction) makes all active detection
channels hopelessly weak.}
\end{finding}


%=============================================================================
\section{Reconciling with the W4i10 Bound: Where $4.2\ee{-10}$ Comes From}
\label{sec:reconcile}
%=============================================================================

\subsection{The W4i10 derivation examined}

Re-reading the W4i10 P2 update (Theorem~2.1, Eq.~(2.6)---(2.7)),
the bound $\lbone < 4.2\ee{-10}$ was derived from the formula:
\begin{equation}
\lbone < \sqrt[4]{\frac{P_{\rm noise}\,\Meff^2\,\Opsi^4}
{P_{\rm emit}\,\Heff^4\,\QEM^4}}
\tag{W4i10 Eq.~(2.7)}
\end{equation}

with $P_{\rm noise} = 10^{-26}$~W, $P_{\rm emit} = 10^6$~W,
$\Meff = 3.01\ee{6}$~kg, $\Heff = 3\ee{-5}$, $\QEM = 10^{11}$,
$\Opsi = 8.17\ee{9}$~rad/s.

Evaluating:
\begin{align}
\frac{P_{\rm noise}\,\Meff^2\,\Opsi^4}{P_{\rm emit}\,\Heff^4\,\QEM^4}
&= \frac{10^{-26}\times(3.01\ee{6})^2\times(8.17\ee{9})^4}
{10^6\times(3\ee{-5})^4\times(10^{11})^4} \nonumber\\
&= \frac{10^{-26}\times 9.06\ee{12}\times 4.46\ee{39}}
{10^6\times 8.1\ee{-19}\times 10^{44}} \nonumber\\
&= \frac{4.04\ee{26}}{8.1\ee{31}} = 4.99\ee{-6}. \label{eq:w4i10_check}
\end{align}

Thus $\lbone < (4.99\ee{-6})^{1/4} = (4.99)^{1/4}\ee{-1.5}
= 1.49\ee{-1.5} = 1.49\times 3.16\ee{-2} = 4.7\ee{-2}$.

\textbf{This gives $\lbone < 0.047$}, not $4.2\ee{-10}$.

\begin{finding}[W4i10 Bound: Arithmetic Error Identified]
\label{find:arithmetic}
\corrected{The W4i10 P2 bound of $\lbone < 4.2\ee{-10}$ contains
a significant error. Tracing through the W4i10 formula with the
stated parameter values gives $\lbone < 0.047$, which is
unconstraining (consistent with Findings~\ref{find:lsw_fails}
and~\ref{find:q_anomaly_fails}).

The likely source of the error: the W4i10 formula uses
$\Heff = 3\ee{-5}$, which is a dimensionless coupling parameter.
If $\Heff$ was intended to represent $4\pi G/c^4 = 9.3\ee{-45}$
(the gravitational-to-EM conversion factor), the numerical
discrepancy would be accounted for. Specifically:
\begin{equation}
\frac{(3\ee{-5})^4}{(9.3\ee{-45})^4} = \frac{8.1\ee{-19}}{7.5\ee{-177}}
= 1.1\ee{158},
\end{equation}
which is far too large.

Alternatively, if $P_{\rm emit} = 10^6$~W is actually the
\emph{circulating} power (not emitted power), and $\Heff$ encodes
the EM-to-psi conversion efficiency including the $G/c^4$ factor,
then the formula might be correct with a different interpretation
of $\Heff$ than we have been able to verify.}

\textbf{Until this discrepancy is resolved, the $4.2\ee{-10}$ bound
and its W4i11 correction to $7.1\ee{-10}$ should be treated as
UNVERIFIED. The SQMS bound of $1.56\ee{-9}$ remains the authoritative
bound.}
\end{finding}


\subsection{The physical barrier}

The fundamental issue is quantitative:
\begin{equation}
\frac{4\pi G}{c^4} = 9.3\ee{-45}~\text{s}^2/(\text{kg}\cdot\text{m}).
\label{eq:G_over_c4}
\end{equation}

This is the factor that converts EM energy density ($\sim 40$~J/m$^3$
in an SRF cavity) to a psi-field source. Even with $Q = 10^{11}$
buildup and $|\lambda-1| \sim 1$, the resulting psi-field perturbation is:
\begin{equation}
\dpsiEM \sim |\lam-1|\,\frac{4\pi G}{c^4}\,u_{\rm EM}\,\frac{Q}{\omega}
\sim 1\times 9.3\ee{-45}\times 40\times\frac{10^{11}}{8.17\ee{9}}
= 4.6\ee{-33}.
\label{eq:dpsi_magnitude}
\end{equation}

The reconversion in the receiver produces a power:
\begin{equation}
P_{\rm recv} \sim \omega\,U_0\,(\dpsiEM)^2 \sim 8.17\ee{9}\times 0.4
\times(4.6\ee{-33})^2 = 6.9\ee{-56}~\text{W}.
\label{eq:P_reconvert}
\end{equation}

Against a noise floor of $10^{-26}$~W, this is $10^{30}$ below
detectability \emph{even for $|\lam-1| = 1$}. This confirms that
the Dark SRF LSW channel cannot constrain $\lbone$ at any
physically interesting level.


%=============================================================================
\section{What the SQMS Bound Actually Constrains}
\label{sec:sqms_bound}
%=============================================================================

The existing SQMS bound of $\lbone < 1.56\ee{-9}$ (W3i3 authoritative)
was derived from a \emph{different} physical channel: the $Q_0$ anomaly
in SRF cavities where the psi-field coupling modifies the surface
impedance. This channel does \emph{not} involve the gravitational
coupling $G/c^4$ directly; instead, it uses the electromagnetic
constitutive relation $n = e^\psi$ with $\psi = \psigrav$ (the
gravitational background) producing a $(\lam-1)$-dependent shift
in the effective electromagnetic properties of the cavity.

This bound was derived in the SQMS context as:
\begin{equation}
\frac{\delta Q_0}{Q_0} = (\lam-1)\,\frac{\delta\psi_{\rm grav}}{d_{\rm skin}}
\,R_{\rm cavity},
\label{eq:sqms_mechanism}
\end{equation}
where $\delta\psi_{\rm grav}/d_{\rm skin}$ is the gradient of the
gravitational psi-field across the skin depth of the superconductor,
and $R_{\rm cavity}$ is a geometric factor.

\textbf{This is a completely different physical mechanism from the
Dark SRF LSW channel.} The SQMS bound uses the \emph{passive}
gravitational background (Earth's potential gradient), not the
\emph{active} EM-sourced psi-perturbation. This is consistent with
the Passive Superiority Theorem (W3i5).


%=============================================================================
%=============================================================================
\part{Multi-Experiment Reinterpretation}
\label{part:multi}
%=============================================================================
%=============================================================================


%=============================================================================
\section{Haloscope Experiments: ADMX, HAYSTAC, ORGAN}
\label{sec:haloscopes}
%=============================================================================

\subsection{Why standard haloscope reinterpretation is different from LSW}

Standard haloscope experiments (ADMX, HAYSTAC, ORGAN) search for
\emph{dark matter} particles that convert to photons in a strong
magnetic field. The signal chain is:
\begin{equation}
\phi_{\rm DM} + B_{\rm ext} \to \gamma_{\rm signal}
\label{eq:haloscope_chain}
\end{equation}
where $\phi_{\rm DM}$ is the dark matter candidate (axion, dark photon,
or scalar) and $B_{\rm ext}$ is the external solenoid field.

For DFD reinterpretation, the question is whether a background $\psi$-field
could mimic this signal. There are two scenarios:

\textbf{Scenario A}: DFD as dark matter. If $\psi$-field fluctuations
constitute a fraction of the dark matter density, they would couple
to the cavity through $u_{\rm EM}$. However, DFD's $\psi$ is a
deterministic classical field, not a quantum field with particle-like
occupation numbers. There is no ``$\psi$ dark matter particle'' to
detect.

\textbf{Scenario B}: Background $\psi$-field as systematic. The
gravitational $\psigrav$ from Earth's potential produces a static
shift in the cavity properties. This does not produce an oscillating
signal at the cavity frequency and therefore does not mimic a dark
matter signal.

\subsection{Form factor comparison: DFD vs dilaton vs dark photon}

\begin{table}[H]
\centering
\caption{Form factor comparison for TM$_{010}$ cylindrical cavity
in a solenoid field.}
\label{tab:formfactor_comparison}
\renewcommand{\arraystretch}{1.3}
\begin{tabular}{@{}lccl@{}}
\toprule
\textbf{Coupling} & \textbf{Source term} & \textbf{$C$ (TM$_{010}$)}
& \textbf{Caveat} \\
\midrule
Axion ($aF\tilde{F}$) & $\mathbf{E}\cdot\mathbf{B}_{\rm ext}$
& 0.69 & Requires $B_{\rm ext}$ \\
Dark photon ($\chi F'F$) & $\mathbf{E}\cdot\hat{\varepsilon}_{A'}$
& 0.50 & Polarization-dependent \\
Dilaton ($\sigma F^2$) & $E^2 - B^2$
& $\sim 10^{-3}$ & Near-cancellation \\
DFD ($\psi\,u_{\rm EM}$) & $E^2 + B^2$
& 0.40 & No cancellation \\
\bottomrule
\end{tabular}
\end{table}

The DFD form factor for haloscope-type experiments is comparable
to the dark photon form factor (0.40 vs 0.50), and dramatically
better than the dilaton (0.40 vs $\sim 10^{-3}$). This is because
DFD couples to \emph{total} energy density, not the Lorentz-invariant
$(E^2 - B^2)$.

However, haloscope experiments cannot be reinterpreted for DFD because
there is no oscillating $\psi$-field signal to detect (the gravitational
background is static). The DFD ``signal'' would be a DC shift in cavity
properties, not an AC signal at the dark matter mass frequency.

\subsection{ORGAN scalar limits and DFD}

McAllister et al.\ (arXiv:2212.01971, Annalen der Physik 2024) computed
scalar/dilaton dark matter limits from ORGAN Phase~1a. Their scalar
form factor is low ($C_\phi \sim 10^{-3}$) because ORGAN uses a uniform
solenoid field and the dilaton couples to $(E^2 - B^2)$.

\textbf{DFD reinterpretation}: If we hypothetically asked ``what if the
DFD $\psi$-field had an oscillating component at the ORGAN frequency
(26.5~GHz)?'', the form factor would be $\sim 0.40$ (using $u_{\rm EM}$
coupling), giving a bound:
\begin{equation}
\lbone_{\rm ORGAN} \sim \lbone_{\rm dilaton}\times
\sqrt{\frac{C_\phi^{\rm dilaton}}{C_\phi^{\rm DFD}}}
\sim 10^{-2}\times\sqrt{\frac{10^{-3}}{0.40}}
\sim 10^{-2}\times 0.05 = 5\ee{-4}.
\label{eq:organ_reinterpret}
\end{equation}

This is not competitive with the SQMS bound.


%=============================================================================
\section{SRF Cavity Experiments Beyond Dark SRF}
\label{sec:other_srf}
%=============================================================================

\subsection{SQMS single-cavity DM search (arXiv:2208.03183)}

Cervantes et al.\ (Phys.\ Rev.\ D~110, 043022, 2024) used a non-tunable
SRF cavity at 1.3~GHz with $Q \sim 10^{10}$ to search for dark photon
dark matter, achieving $\chi < 1.5\ee{-16}$ at $m_{A'} = 5.35~\mu$eV.

\textbf{DFD reinterpretation}: This is a DM haloscope search, not LSW.
The same argument applies: there is no oscillating $\psi$-field at
$5.35~\mu$eV to detect in DFD. No useful bound on $\lbone$ is obtained.

\subsection{Dark SRF Phase~2 (planned)}

The next generation Dark SRF aims for 2.6~GHz with $Q > 10^{11}$ and
improved shielding. The LSW sensitivity would improve by $\sim 10\times$
in $\chi$. As shown in Section~\ref{sec:reconcile}, this still does not
produce a useful DFD bound through the LSW channel.

\subsection{ADMX and HAYSTAC}

ADMX (1--2~GHz range) achieved $g_{a\gamma\gamma}$ sensitivity at the
DFSZ level for axions. HAYSTAC (4--5~$\mu$eV range) used squeezed-state
receivers. Neither can be reinterpreted for DFD (haloscope, not LSW;
static $\psi$ background, no oscillating signal).

\begin{finding}[Multi-Experiment Reinterpretation: Limited Applicability]
\label{find:multi_experiment}
\honest{No existing haloscope experiment (ADMX, HAYSTAC, ORGAN,
SQMS DM search) can be usefully reinterpreted for DFD bounds on
$\lbone$. The fundamental reason: DFD's $\psi$-field is a classical
gravitational background, not a dark matter particle. Haloscopes
search for oscillating signals at the dark matter mass frequency;
the DFD signal is a static (or very slowly varying) shift in cavity
properties.

The only SRF experiment pathway to a DFD bound is the passive
$Q_0$-measurement approach (the SQMS bound), which constrains how
$\lbone$ modifies the surface impedance through the gravitational
$\psi_{\rm grav}$ background. This is the mechanism behind the
existing $\lbone < 1.56\ee{-9}$ bound (W3i3).}
\end{finding}


%=============================================================================
%=============================================================================
\part{Frequency and Mass Dependence}
\label{part:frequency}
%=============================================================================
%=============================================================================


%=============================================================================
\section{Mass Dependence of the Passive $Q_0$ Bound}
\label{sec:mass_dependence}
%=============================================================================

\subsection{The DFD psi-field mass}

In DFD, the $\psi$-field is effectively massless in the linearized
regime ($|\nabla\psi|/a_* \gg 1$, Newtonian limit). In the MOND regime
($|\nabla\psi|/a_* \ll 1$), the nonlinear kinetic term creates an
effective mass-like term through self-interaction, but this is relevant
only at galactic scales ($a < a_0$).

For laboratory experiments, $\psi$ is effectively massless:
$m_\psi \approx 0$. The Compton wavelength is infinite.

\subsection{Frequency dependence of the passive $Q_0$ bound}

The SQMS $Q_0$ bound depends on the cavity frequency through:
\begin{equation}
\lbone \propto \frac{1}{\sqrt{Q_0(f)\,u_{\rm EM}(f)\,f^{-1}}}.
\label{eq:freq_scaling}
\end{equation}

The key frequency dependences:
\begin{itemize}[nosep]
\item $Q_0(f)$: SRF quality factors decrease at higher frequencies
  due to surface resistance scaling as $R_s \propto f^2$ (BCS theory).
  At 1.3~GHz: $Q_0 \sim 10^{11}$. At 10~GHz: $Q_0 \sim 10^8$--$10^9$.
\item $u_{\rm EM}(f)$: Stored energy density is geometry-dependent
  but generally lower in smaller (higher-frequency) cavities.
\item The $\psi$-field coupling does not depend on frequency
  (it couples to $u_{\rm EM}$, which is frequency-independent
  at fixed stored energy).
\end{itemize}

\subsection{Optimal frequency for the passive bound}

The passive $Q_0$ bound tightens with:
\begin{equation}
\text{Figure of merit} = Q_0(f)\,\sqrt{U_0(f)},
\label{eq:FoM}
\end{equation}
which is maximized at the lowest practical SRF frequency where
$Q_0$ is highest. The 1.3~GHz TESLA frequency is near-optimal
for Nb SRF technology.

Going to higher frequencies (e.g., 10~GHz for the frequency-optimized
P4 conversion at 32~GHz) would \emph{weaken} the passive bound due
to the $Q_0 \propto f^{-2}$ scaling of BCS losses.

\subsection{Mass range where the bound applies}

Since the DFD $\psi$-field is massless, the passive $Q_0$ bound
applies \emph{at all cavity frequencies simultaneously}. There is
no mass-resonance condition to satisfy. The bound is:
\begin{equation}
\lbone < 1.56\ee{-9} \qquad\text{for all $m_\psi$ (including $m_\psi = 0$)}.
\label{eq:universal_bound}
\end{equation}

This is in contrast to dark photon bounds, which apply only in a
narrow mass window around $m_{A'} = \hbar\omega_{\rm cav}/c^2$.

\begin{finding}[Frequency/Mass Dependence]
\label{find:frequency}
\confirmed{The DFD $\psi$-field passive $Q_0$ bound is:
\begin{enumerate}[nosep]
\item \textbf{Mass-independent}: applies for all $m_\psi$, including
  $m_\psi = 0$.
\item \textbf{Frequency-optimized at 1.3~GHz}: the BCS $R_s \propto f^2$
  scaling means lower frequencies give better $Q_0$ and tighter bounds.
\item \textbf{Not improvable by frequency tuning}: unlike haloscope
  experiments, there is no resonance peak to find.
\end{enumerate}

The only path to tightening the bound is increasing $Q_0$ at fixed
frequency (the SQMS Phase~I approach) or using multiple cavities for
cross-correlation (the Phase~I 4-cavity array).}
\end{finding}


%=============================================================================
%=============================================================================
\part{Conclusions and Programme Impact}
\label{part:conclusions}
%=============================================================================
%=============================================================================


%=============================================================================
\section{Corrected Detection Landscape}
\label{sec:corrected_landscape}
%=============================================================================

\begin{finding}[Programme-Level Assessment]
\label{find:programme}
This iteration identifies a significant issue with the W4i10--W4i11
Dark SRF reinterpretation chain:

\begin{enumerate}
\item \textbf{The $\lbone < 4.2\ee{-10}$ (W4i10) and $\lbone < 7.1\ee{-10}$
(W4i11) bounds are UNVERIFIED.} Independent recomputation through three
translation paths (LSW power match, Q-anomaly, dimensionless coupling)
all yield $\lbone \gg 1$ from the Dark SRF data. The gravitational
coupling $G/c^4 = 9.3\ee{-45}$ makes the active EM-to-psi conversion
hopelessly weak, consistent with the Passive Superiority Theorem (W3i5)
and the Volume Cancellation Theorem (W3i5).

\item \textbf{The SQMS bound of $\lbone < 1.56\ee{-9}$ remains
authoritative.} This bound uses the passive $Q_0$ mechanism (gravitational
background $\psi_{\rm grav}$), not the active LSW mechanism. The passive
approach avoids the $G/c^4$ suppression by leveraging the \emph{existing}
gravitational potential gradient, not an EM-sourced perturbation.

\item \textbf{No haloscope experiment can be usefully reinterpreted
for DFD.} The DFD $\psi$-field is not a dark matter particle. Haloscopes
search for oscillating signals at the DM mass frequency; DFD produces
static shifts.

\item \textbf{The DFD coupling dictionary is now established.} The
key distinction from dilaton coupling ($E^2 - B^2$, which cancels in
cavities) is that DFD couples to $(E^2 + B^2)$, giving form factors
comparable to dark photon coupling ($\sim 0.4$ vs $\sim 0.5$).

\item \textbf{The SQMS Phase~I programme remains the correct
detection pathway.} The $Q_0$-ratio diagnostic (DFD: 0.49 vs BCS: 3.41)
is independent of the absolute bound and provides a decisive test.
\end{enumerate}
\end{finding}


%=============================================================================
\section{Updated Detection Ranking}
\label{sec:ranking}
%=============================================================================

\begin{table}[H]
\centering
\caption{Corrected detection ranking (W4i12). Dark SRF reinterpretation
DEMOTED.}
\label{tab:corrected_ranking}
\renewcommand{\arraystretch}{1.3}
\begin{tabular}{@{}clccc@{}}
\toprule
\textbf{Rank} & \textbf{Experiment} & \textbf{Cost}
& \textbf{Reach $\lbone$} & \textbf{Status} \\
\midrule
\rowcolor{rowhi}
1 & SQMS Phase~I (4-cavity) & \$3.2M/24mo & $7.8\ee{-10}$ & Correct \\
\rowcolor{rowhi}
2 & LIGO O4 archival 6th channel & \$0.1--0.5M & $\sim 1.5\ee{-14}$ & Correct \\
\rowcolor{rowhi}
3 & ESS Channel~B & \$2--5M & $\sim 8\ee{-13}$ & Correct \\
\rowcolor{rowhi}
4 & GPS 59~ps archival & \$50--100K & $\lam$-dependent & Correct \\
\rowcolor{rowhi}
5 & Si$_3$N$_4$ MEMS in-gap & \$200--500K & $Q_\psi$ threshold & Correct \\
\rowcolor{rowlo}
--- & Dark SRF reinterpretation & \$0 & UNVERIFIED & \textbf{DEMOTED} \\
\bottomrule
\end{tabular}
\end{table}

The authoritative bound remains $\lbone < 1.56\ee{-9}$ (SQMS, W3i3).


%=============================================================================
\section{Recommendations}
\label{sec:recommendations}
%=============================================================================

\begin{enumerate}
\item \textbf{URGENT}: The W4i10 P2 derivation of $\lbone < 4.2\ee{-10}$
  from Dark SRF data needs a careful line-by-line audit. The parameter
  $H_n = 3\ee{-5}$ in Theorem~2.1 needs to be traced to its dimensional
  origin. If it encodes $G/c^4$ (gravitational coupling), the bound is
  hopelessly weak. If it encodes a different physical mechanism (e.g.,
  the passive $\psi_{\rm grav}$ background modifying $Q_0$), the
  derivation path needs to be made explicit.

\item \textbf{The coupling dictionary (Table~\ref{tab:dictionary}) should
  be incorporated into the master corpus.} The key result --- DFD couples
  to $(E^2 + B^2)$, not $(E^2 - B^2)$ like dilaton --- has implications
  for all future cavity experiments.

\item \textbf{SQMS Phase~I remains the gateway experiment.} The Q-ratio
  diagnostic is robust against all bound uncertainties.

\item \textbf{A dedicated study of whether the passive $Q_0$ mechanism
  can be applied to Dark SRF data is warranted.} If the Dark SRF $Q_0$
  measurements show any anomalous frequency or temperature dependence
  beyond BCS, this could constrain $\lbone$ through the same mechanism
  as the SQMS bound --- but with better $Q_0$ ($10^{11}$ vs $5\ee{10}$).
\end{enumerate}


%=============================================================================
\section{Sources}
\label{sec:sources}
%=============================================================================

\begin{itemize}[nosep]
\item Dark SRF PRL 130, 261801 (2023):
  \url{https://journals.aps.org/prl/abstract/10.1103/PhysRevLett.130.261801}
\item Dark SRF improved: arXiv:2510.02427,
  \url{https://arxiv.org/abs/2510.02427}
\item SQMS DM search: arXiv:2208.03183, Phys.\ Rev.\ D~110, 043022 (2024),
  \url{https://arxiv.org/abs/2208.03183}
\item ORGAN scalars: arXiv:2212.01971, McAllister et al., Annalen der Physik 2024,
  \url{https://arxiv.org/abs/2212.01971}
\item Flambaum dilaton: arXiv:0711.2858,
  \url{https://ar5iv.labs.arxiv.org/html/0711.2858}
\item Flambaum scalar cavity: arXiv:2207.14437, Phys.\ Rev.\ D~106, 055037 (2022),
  \url{https://arxiv.org/abs/2207.14437}
\item LSW cavities: arXiv:2303.15996, arXiv:2402.09899,
  \url{https://arxiv.org/html/2303.15996}, \url{https://arxiv.org/html/2402.09899}
\item Berlin et al.\ 2020: FERMILAB-PUB-20-256-T,
  \url{https://lss.fnal.gov/archive/2020/pub/fermilab-pub-20-256-t.pdf}
\item Fermilab News (2023):
  \url{https://news.fnal.gov/2023/07/dark-srf-experiment-at-fermilab-demonstrates-ultra-sensitivity-for-dark-photon-searches/}
\item ADMX 2025: arXiv:2504.07279,
  \url{https://arxiv.org/html/2504.07279}
\item HAYSTAC Phase II: PRL 134, 151006 (2025),
  \url{https://link.aps.org/doi/10.1103/PhysRevLett.134.151006}
\end{itemize}

\end{document}
